% !PRINTED 393
\weberpage{393}
Dieser Fall führt also immer auf eine reelle Doppeltangente.

Wenn aber \gesperrt{zweitens} das Complexpaar 1), 2) azygetisch ist, so
enthält jedes Paar des einen Complexes ein Element, was auch im
anderen Complexe vorkommt, und ein Element, was im anderen nicht
vorkommt. Kommt also etwa $\eta$ im Complexe 2) vor, so können wir,
ohne Beeinträchtigung der Allgemeinheit $\eta = \eta_1'$, $\eta' = \eta_1$
annehmen, und erhalten folgende Paare in 1) und 2)
\weq{2}{
\begin{array}{rl}
1) & \xi\eta,\ \xi_1\eta' \\
2) & \xi_1'\eta,\ \xi'\eta'.
\end{array}
}
Wenn ferner $\eta_2$ in 1) und 2) vorkommt, so kann $\eta_2$ nicht gleich
$\xi_2'$ sein, weil sonst 1) das Paar $\xi_2\xi_2'$ enthielte und folglich
reell wäre.

Wenn $\eta_2 = \eta_2'$ ist, so ist $\eta_2$ \gesperrt{reell}. Ist aber $\eta_2$ gleich
einem Element der drei letzten Paare von 2), so können wir es ohne
Beschränkung $= \eta_3'$ annehmen, also $\eta_2 = \eta_3'$, $\eta_3 = \eta_2'$.

Dann sind aber nach dem Satze \S.\,96, 3. sowohl $\eta, \eta_2, \eta_2'$ als
$\eta, \eta', \eta_2$ azygetisch. Die beiden Paare $\eta\eta_2$, $\eta'\eta_2'$
bestimmen zwei Complexe, deren zweiter weder $\eta$ noch $\eta_2$ enthält,
und die daher ein \gesperrt{syzygetisches Paar} bilden. Damit sind wir auf
die erste Annahme zurückgeführt und unser Satz 1. ist bewiesen.

Es giebt also immer mindestens \gesperrt{zwei reelle} Doppeltangenten.

\gesperrt{2. Es giebt immer mindestens ein System von vier reellen
syzygetischen Doppeltangenten.}

Beim Beweise dieses Satzes gehen wir aus von einem nach 1. immer
existirenden reellen Paare $x\,y$, und betrachten den reellen Complex,
der durch dies Paar bestimmt ist. Da die conjugirten Doppelpaare je
zwei Paare sind, so muss unter den sechs Paaren dieses Complexes
entweder ein zweites reelles Paar vorkommen, und dann trifft der Satz
2. zu, oder er muss ein conjugirtes Paar $\xi\xi'$ enthalten.

Im letzteren Falle betrachten wir das durch diese beiden Paare
bestimmte syzygetische Complextripel, in dem alle 28 Doppeltangenten
vorkommen (\S.\,97, 4.):
\[
\begin{array}{rl}
1) & x\,y,\ \xi\xi' \\
2) & x\xi,\ y\xi',\ \xi_1\eta_1 \\
3) & x\xi',\ y\xi,\ \xi_1'\eta_1',
\end{array}
\]
%CONT>>
